\chapter{Conclusion}\label{chap:conclusion}

This report addressed the three connected tasks of the project \textit{Towards
Robust Forming Limit Curve Identification: Experiments and Modeling}. Task~I
generated the experimental forming-limit data used to study material
orientation, sheet thickness, punch diameter, and evaluation-window effects.
Task~II extended the identification workflow by examining a Min--Stoughton
curvature-based criterion. Task~III supplied an automated numerical framework
for generating controlled simulation data for the same specimen and punch
configurations.

The experimental campaign showed that the RD and TD forming limit curves of the
AA5754 sheets are indistinguishable within the repeatability of the available
tests. The punch-diameter study also showed that applying the standard
ISO~12004-2 window unchanged to non-standard punch diameters can distort the
extracted limit, especially for the \SI{40}{\milli\meter} punch. A
punch-scaled fitting window therefore appears necessary before the size-effect
campaign can be interpreted quantitatively.

The Min--Stoughton study demonstrated that the curvature workflow is
operational on laboratory DIC exports and can produce onset frames in the range
of the ISO~12004-2 and Volk--Hora results. However, the sensitivity to matrix
length, reference-frame choice, and DIC coverage means that the method is not
yet calibrated for production FLC extraction.

On the numerical side, the main implemented result is a reproducible and
extensible FEA pipeline. The model-generation step creates the tooling,
specimen, contact definitions, material assignment, output requests, and
boundary conditions automatically. Solver jobs are submitted through
\software{SLURM}, and the post-processing stage extracts force--displacement
histories, energy histories, strain paths, forming-limit points, and diagnostic
media from the \software{Abaqus} output database.

In summary, the report connects the experimental, data-evaluation, and
simulation parts of the project into one workflow. The remaining work is mainly
quantitative: completing the numerical sensitivity studies, validating the
simulation output against the experimental FLCs, and refining the identification
criteria once larger datasets are available.
